\begin{table}[tp]
\centering
\caption{Power-geometry acceleration with relative gradient error. \href{https://github.com/k3nfalt/longtails/blob/main/LongTails/Setups/PowerGeometryRelativeErrorConvex.lean}{PowerGeometryRelativeErrorConvex}.}\label{tab:PowerGeometryRelativeErrorConvex}
\parbox{\linewidth}{\footnotesize In a finite-dimensional real normed space, $f$ is differentiable and convex, $\|Df(x)-Df(y)\|_*\le L\|x-y\|$, and a global minimizer satisfies $\|x^*-x_0\|\le R$. Here $L,c>0$, $p>2$, $R\ge0$, and $0\le\alpha\le1/4$. The differentiable function $\psi(x)=\|x-x_0\|^2/2$ satisfies $D_\psi(y,x)\ge c\|y-x\|^p/(\|y-x_0\|+\|x-x_0\|)^{p-2}$ for all $x,y$ (zero at $x=y=x_0$). Every oracle reply obeys $\|g-Df(x)\|_*\le\alpha\|Df(x)\|_*$, allowing biased, call-dependent errors. Target: $f(x_K)-f(x^*)\le\varepsilon$, $\varepsilon>0$, almost surely under every execution law; the constructed squared-norm dual-averaging output satisfies it for every seed.}\par\smallskip
\begin{threeparttable}
\footnotesize
\setlength{\tabcolsep}{4pt}\renewcommand{\arraystretch}{1.25}
\begin{tabularx}{\linewidth}{@{}>{\raggedright\arraybackslash}p{0.20\linewidth}>{\raggedright\arraybackslash}p{0.32\linewidth}>{\raggedright\arraybackslash}p{0.25\linewidth}>{\raggedright\arraybackslash}X@{}}
\toprule
\textbf{Result} & \textbf{Oracle complexity} & \textbf{Novelty} & \textbf{Limitations} \\
\midrule
\textbf{(1)} \href{https://github.com/k3nfalt/longtails/blob/main/LongTails/Theorems/PowerGeometryRelativeErrorAcceleration.lean}{Power\allowbreak Geometry\allowbreak Relative\allowbreak Error\allowbreak Acceleration} \textbf{(Ours)} $(\star\star\star)$ \href{https://github.com/k3nfalt/longtails/blob/main/scripts/reviews/PowerGeometryRelativeErrorAcceleration.md}{(review)}
& $\mathcal{O}\!\left(1+(CT)^{p/(p+2)}+\alpha^2T\right)$ gradient calls.
& & \\
\midrule[1.2pt]
\textbf{(2)} \citep{pgDAspremont2016v3}, Cor.~4.10, \S5.1.2
& $\mathcal{O}_p\!\left(1+T^{p/(p+2)}\right)$ exact-gradient calls. Special case: $\ell_p^d$, $\alpha=0$, optimization on $B_p(x_0,R)$ containing a global minimizer.
& \textbf{(1)}: extends to biased relative error; matches the exact-gradient rate for fixed $p$. Updates need neither $R$, $\varepsilon$, nor horizon; the stopping budget still does.
& \textbf{(1)} has larger $p$-dependent constants; on $\ell_p^d$, $C=p(p-1)2^{7(p-2)/2}$ (1536 at $p=4$). \\
\midrule
\textbf{(3)} Derived bound: \citep{pgVasinGasnikov2026v1}, Thm.~7.2
& $\mathcal{O}\!\bigl(1+(d^s\sqrt T+\alpha d^{3s}T)$\par ${}\cdot[1+\log(d^{2s}T)]\bigr)$ gradient calls. On $\ell_p^d$ via Euclidean conversion, with $d^{2s}T>1$ and $d^s\alpha\le1/3$.
& \textbf{(1)}: dimension-free; for fixed $p$ and $d\gtrsim T^{p/(p+2)}$, removes dimension factors and logarithm, with $\alpha^2T$ instead of $\alpha d^{3s}T$. At $p=4$, $d=T^2$, $\alpha=T^{-1/6}\le1/4$: $\mathcal{O}(1+T^{2/3})$; \textbf{(3)} is inapplicable.
& \textbf{(1)} loses at $\alpha=0$, fixed $p,d$, $T\to\infty$: $T^{p/(p+2)}$ versus $\sqrt T(1+\log T)$. At $d=1$, excludes $1/4<\alpha\le1/3$. \\
\midrule
\textbf{(4)} SODA \citep{pgPethick2026v1}, Cor.~3.2
& $\mathcal{O}_{p,\tau_q}\!\bigl(1+T^{p/(p+2)}+(\sigma_qR/\varepsilon)^p\bigr)$ stochastic-gradient calls, one per step. On $\ell_p^d$: unbiased gradients with $q$-moment gradient variation (Assumptions~3--4); expected objective gap $\le\varepsilon$.
& \textbf{(1)}: matches the rate on the exact-gradient overlap; removes $R$ from updates. Covers arbitrary relative bias with pathwise accuracy; no stochastic-rate comparison outside that overlap.
& \textbf{(1)} does not cover additive heavy-tailed noise allowed by \textbf{(4)}; the noisy oracle models are incomparable. \\
\bottomrule
\end{tabularx}
\begin{tablenotes}[flushleft]\footnotesize
\item[] \textbf{Abbr:} $T=LR^2/\varepsilon$; $C=\max\{1,32^{(p-2)/2}/c\}$; $D_\psi(y,x)=\psi(y)-\psi(x)-D\psi(x)[y-x]$; $\|\cdot\|_*$ is the dual norm. For $\ell_p^d$, $d$ is dimension, $s=1/2-1/p$, and $c=2^{2-p}/[p(p-1)]$ is admissible. $\mathcal{O}_p$ permits constants depending on fixed $p$. In \textbf{(4)}, $q=p/(p-1)$; $\sigma_q$ and $\tau_q\ge1$ are the additive scale and multiplicative factor in the $L^q$ gradient-variation bound; $\mathcal{O}_{p,\tau_q}$ permits dependence on these subscripts.
\end{tablenotes}
\end{threeparttable}
\end{table}
\clearpage
